\chapter{Appendix 01: Alpha Fabrication}

\section{The Legacy Paradigm \\\& Its Failures}
Legacy fabrication relied on proprietary CAD software tethered to cloud DRM and centralized processing. If the internet went down, production stopped. Interfaces were hyper-complex, demanding intense cognitive load, and ignored the physical reality and thermodynamic cost of the workshop environment.

\section{The Incremental Automation Pathway}
Phase 1 transitions from cloud CAD to local-first CRDT-based part repositories. Phase 2 introduces spatial UI, projecting assembly instructions directly onto the physical materials via AR. Phase 3 involves the machines automatically fetching thermodynamic rendering limits and adjusting their operations, requiring human interaction only for safety overrides.

\section{The Human Boundary (Limits of Automation)}
Automation halts at the physical safety boundary. Before a high-energy fabrication tool (e.g., a laser cutter or CNC mill) is engaged, the interface strips away all rich graphics. A low-power, high-contrast prompt forces the operator to physically push an approval button, ensuring explicit human oversight.

\section{Interface Mechanics (The "How")}
Spatial interfaces overlay toolpath projections onto raw materials. CRDTs ensure multiple operators can edit a design simultaneously without internet access. Yet, a missing mechanism is identified: real-time haptic feedback interfaces to warn operators of physical tool deviations without relying on visual AR.
